\chapter{Experiments that could test the connections}\label{ch:experimental-bridge}

\section{A theorem suggests a question, not a result}
The hierarchy gives several exact predictions under declared assumptions. Those predictions can guide a future experimental programme, but they do not choose the apparatus, intervention sizes, tolerances or sample counts. Those decisions require a separate prospective design and authorization.

The most useful sequence begins with comparisons whose interpretation is sharp. Establish a pair effect under a controlled quadratic field, test a triple null with certified additive response, and only then seek a positive higher-order benchmark. Here ``positive benchmark'' means a known nonzero target that the method should detect, not necessarily a positive EBU sign.

\section{First establish the pair correction}
A pair study would compare all four required endpoints under one independently calibrated $H$ and fixed increments $h_A,h_B$. The prediction is
\begin{equation}
 m_E(AB)=-h_A^{\mathsf T}Hh_B.
\end{equation}
The design should include conditions under which the correction is nonzero and, where appropriate, a structurally justified zero case. It should verify accepted physical increments rather than assume that identical control settings produce identical changes in every context.

This study could reveal a baseline error, an omitted shared coordinate or a response dependence before the analysis becomes more complicated. Its societal relevance is foundational: a future system must first show it can avoid double-counting a simple shared effect before claiming to value large cooperative projects.

\section{Then test the additive quadratic triple null}
With three fixed increments and a quadratic field, all eight alternatives predict $m_E(ABC)=0$. A valid physical null requires a complete admissible cube, common field and baseline, and a covariance model that includes shared calibration and measurements.

A nonzero result would reject the joint null assumptions at the stated uncertainty, not prove one selected explanation. A zero result would support compatibility within the tested resolution and domain, not universal absence of higher-order interaction. Controls should help separate response nonlinearity, field error and measurement drift.

The null is valuable because it can fail. If every nonzero coefficient is removed as a bad example or absorbed into a newly fitted field after observation, the experiment no longer tests the original proposition.

\section{Introduce a controlled cubic target}
A physically justified local cubic contribution can supply the exact algebraic target $m_E(ABC)=-\lambda abc$ under additive increments. A real design must specify its valid domain, stability, boundary and calibration; the unbounded scalar teaching cubic is not itself an apparatus proposal.

The purpose is to show that the measurement and analysis can detect a known higher-order contrast when it should be present. This complements the quadratic null. Passing only a null can mean a method is insensitive; passing only a nonzero case can conceal bias. The two questions belong together.

\section{Hold the quadratic and change the response}
A separate test could retain an independently calibrated quadratic potential while introducing a characterized nonlinear response. The target must come from that response law, not from the assumption that all nonlinear devices produce a triple.

The constructed map $x=z_1+z_2+z_3+\alpha z_1z_2$ offers an algebraic example with $m_E(123)=-\alpha$. A physical implementation would need its own state, controls and measurements. The experiment should verify the response independently enough to distinguish it from an incorrect potential fit.

This comparison could teach a future EBU system where a dependency resides. Engineering a more additive response and changing a valuation model are different interventions. Knowing which one explains a measured effect can prevent unnecessary changes to an otherwise appropriate field.

\section{Compare two independent branches at equilibrium}
The probability bridge suggests two readings of the same coefficient. Branch A uses physically calibrated potential differences. Branch B estimates the alternating log-density contrast in a declared reference measure. Under the canonical normalization they should agree.

The independence of the branches is essential. If Branch A defines its potential by taking the negative logarithm of Branch B's estimated density, agreement is construction, not validation. Shared measurements or calibration may still create covariance, which must be included in the comparison.

The study also needs compatible endpoint support, positive density where logs are used, density-of-states treatment, and a plan for uncertainty and estimator bias. A thermal entropy interpretation requires the additional P4 conditions; density agreement alone is not a test of reversibility or absence of work.

\section{Cross-field interaction is a later question}
Only after a justified denomination bridge could one compare higher-order contrasts across different fields as quantities on a common scale. Actions must also be matched in physical meaning. Equal labels do not guarantee equal interventions, and mathematically dimensionless potentials need not have interchangeable units of significance.

This leaves an ambitious research direction open without treating it as established. If comparable accounts across domains become defensible, they could help institutions examine tradeoffs that price records alone obscure. If comparability fails, separate vector accounts may remain more informative than an unjustified scalar sum.

\section{Relation to E1a}
The current E1a physical-probabilistic work provides a nearby benchmark and an explicit separation between physical and statistical branches. Its present scientific and execution status remains governed by the frozen repository records. This chapter neither changes its target nor restarts the paused validation track.

The interaction studies described here are proposals beyond the current book integration. No seed, threshold, horizon, instrument or campaign is adopted in these pages. A later design should identify which prerequisites E1a actually establishes and which the new interaction question still needs.

\section{Beyond one favorable coefficient}
An institution aimed at sustaining life support needs evidence about service, reserves, resilience, distribution and access over time. A positive interaction coefficient is one possible piece of information, not a complete success metric. A group can have a favorable interaction while failing to deliver a necessary service or shifting harm outside the boundary.

A long-run study must therefore define its physical domain and comparable alternatives, distinguish natural forcing from actor action, and retain refusal, cycling, adverse and inconclusive outcomes. Boundedness of a continuous potential on a finite closed stock domain is already guaranteed by the domain; it cannot discriminate the success of one actor mechanism.

Similarly, a potential minimum does not prove attraction. A finite action menu can make a compatible reference unreachable. If every allowed transfer changes stocks by an even integer, parity can prevent reaching a reference of the opposite parity. No amount of exact valuation repairs that physical opportunity set.

\section{A credible route to social benefit}
The proposed sequence has a practical purpose: establish that the account measures what it claims, then examine whether people and institutions can use that information well. Small exact benchmarks can reduce uncertainty about the mathematical and physical links. Later studies can ask whether maintenance is better supported, complementary services are coordinated, and essential access is protected.

A failed bridge or adverse institutional result would be useful evidence. It would identify a limitation rather than invalidate every earlier identity. Conversely, a successful benchmark would not authorize broad adoption without further work. The ambition is substantial enough to deserve this sequence of distinct questions.

The final chapters return to the system as a whole: how these pieces might fit into an inspectable institution, and what would count as success beyond a correctly closing account.
